% GENERATED FILE. Edit canonical lessons, then run npm run generate:books.
% canonical-source-hash: fnv1a64:1c9ccb83dc8ecd21
% canonical-lessons: PT-C320-costume, PT-C320-ritual, PT-C320-cerimonia, PT-C320-peregrino, PT-C320-peregrinacao

\chapter{Ritos e costumes}
\label{ch:pt-ritos-e-costumes}
\hlchaptermodality{320}

\begin{chapteropening}
\textbf{By the end of this chapter:} \emph{I can talk about customs, rituals, ceremonies and pilgrimages.}

Complete the last lesson of chapter 320: I can talk about customs, rituals, ceremonies and pilgrimages.
\end{chapteropening}

\section[o costume]{o costume — the custom, the habit}

\label{lesson:PT-C320-costume}



\begin{quote}
Before the new one: say the Portuguese for the sardine feast, then the Portuguese for the street dance.
\end{quote}

\subsection*{You'll want to know: o costume}
\textbf{o costume} — \textquotedblleft{}the custom, the habit\textquotedblright{}.

\begin{quote}
É costume jantar tarde nesta região.
\end{quote}

\begin{quote}
\emph{It is the custom to eat late in this region.}
\end{quote}

\emph{Como de costume} — as usual. Not a party outfit.

\begin{cousinweb}
\emph{costume} comes down from Old Portuguese (EtymDB) and goes back to Latin \emph{cōnsuētūdō} \textquotedblleft{}habit, custom\textquotedblright{}, from \emph{cōnsuēscere} \textquotedblleft{}to get used to\textquotedblright{} (Lewis \& Short).
\end{cousinweb}

\begin{grammarlens}[title={What you've built}]
One more word for telling, explaining and describing at B1.
\end{grammarlens}

\subsection*{Your turn}
\begin{itemize}
\raggedright
  \item \emph{Say it:} \emph{o costume}
  \item \emph{Say it:} \emph{o costume}, once more
  \item \emph{Say it:} a sentence of your own with \emph{o costume}
  \item \emph{Recall:} say the Portuguese for the sardine feast, then the Portuguese for the street dance, then say \emph{o costume} again
\end{itemize}

\begin{tcolorbox}[breakable,colback=teal!4,colframe=teal!35!black,title={Before you move on}]
What does o costume mean? (\textquotedblleft{}The custom, the habit\textquotedblright{}.) Say it once more.
\end{tcolorbox}
\section[o ritual]{o ritual — the ritual}

\label{lesson:PT-C320-ritual}



\begin{quote}
Before the new one: say the Portuguese for the street dance, then the Portuguese for the custom.
\end{quote}

\subsection*{You'll want to know: o ritual}
\textbf{o ritual} — \textquotedblleft{}the ritual\textquotedblright{}.

\begin{quote}
Tomar café depois do almoço é um ritual para ele.
\end{quote}

\begin{quote}
\emph{Having coffee after lunch is a ritual for him.}
\end{quote}

Also an adjective: \emph{uma dança ritual}.

\begin{cousinweb}
\emph{ritual} is a learned borrowing from Latin \emph{rītuālis} \textquotedblleft{}of rites\textquotedblright{}, from \emph{rītus} \textquotedblleft{}a rite, a custom\textquotedblright{} (Lewis \& Short).
\end{cousinweb}

\begin{grammarlens}[title={What you've built}]
One more word for telling, explaining and describing at B1.
\end{grammarlens}

\subsection*{Your turn}
\begin{itemize}
\raggedright
  \item \emph{Say it:} \emph{o ritual}
  \item \emph{Say it:} \emph{o ritual}, once more
  \item \emph{Say it:} a sentence of your own with \emph{o ritual}
  \item \emph{Recall:} say the Portuguese for the street dance, then the Portuguese for the custom, then say \emph{o ritual} again
\end{itemize}

\begin{tcolorbox}[breakable,colback=teal!4,colframe=teal!35!black,title={Before you move on}]
What does o ritual mean? (\textquotedblleft{}The ritual\textquotedblright{}.) Say it once more.
\end{tcolorbox}
\section[a cerimónia]{a cerimónia — the ceremony}

\label{lesson:PT-C320-cerimonia}



\begin{quote}
Before the new one: say the Portuguese for the custom, then the Portuguese for the ritual.
\end{quote}

\subsection*{You'll want to know: a cerimónia}
\textbf{a cerimónia} — \textquotedblleft{}the ceremony\textquotedblright{}.

\begin{quote}
A cerimónia durou uma hora.
\end{quote}

\begin{quote}
\emph{The ceremony lasted an hour.}
\end{quote}

In Brazil: \emph{cerimônia}. \emph{Sem cerimónia} — without fuss.

\begin{cousinweb}
\emph{cerimónia} is a learned borrowing from Latin \emph{caerimōnia} \textquotedblleft{}a sacred rite, reverence\textquotedblright{} (EtymDB).
\end{cousinweb}

\begin{grammarlens}[title={What you've built}]
One more word for telling, explaining and describing at B1.
\end{grammarlens}

\subsection*{Your turn}
\begin{itemize}
\raggedright
  \item \emph{Say it:} \emph{a cerimónia}
  \item \emph{Say it:} \emph{a cerimónia}, once more
  \item \emph{Say it:} a sentence of your own with \emph{a cerimónia}
  \item \emph{Recall:} say the Portuguese for the custom, then the Portuguese for the ritual, then say \emph{a cerimónia} again
\end{itemize}

\begin{tcolorbox}[breakable,colback=teal!4,colframe=teal!35!black,title={Before you move on}]
What does a cerimónia mean? (\textquotedblleft{}The ceremony\textquotedblright{}.) Say it once more.
\end{tcolorbox}
\section[o peregrino]{o peregrino — the pilgrim}

\label{lesson:PT-C320-peregrino}



\begin{quote}
Before the new one: say the Portuguese for the ritual, then the Portuguese for the ceremony.
\end{quote}

\subsection*{You'll want to know: o peregrino}
\textbf{o peregrino} — \textquotedblleft{}the pilgrim\textquotedblright{}.

\begin{quote}
O peregrino caminhou duzentos quilómetros.
\end{quote}

\begin{quote}
\emph{The pilgrim walked two hundred kilometres.}
\end{quote}

A woman is \emph{a peregrina}.

\begin{cousinweb}
\emph{peregrino} comes from Latin \emph{peregrīnus} \textquotedblleft{}foreign, a foreigner\textquotedblright{} (EtymDB), built on \emph{per} \textquotedblleft{}through\textquotedblright{} and \emph{ager} \textquotedblleft{}a field\textquotedblright{}: one who comes from beyond the fields.
\end{cousinweb}

\begin{grammarlens}[title={What you've built}]
One more word for telling, explaining and describing at B1.
\end{grammarlens}

\subsection*{Your turn}
\begin{itemize}
\raggedright
  \item \emph{Say it:} \emph{o peregrino}
  \item \emph{Say it:} \emph{o peregrino}, once more
  \item \emph{Say it:} a sentence of your own with \emph{o peregrino}
  \item \emph{Recall:} say the Portuguese for the ritual, then the Portuguese for the ceremony, then say \emph{o peregrino} again
\end{itemize}

\begin{tcolorbox}[breakable,colback=teal!4,colframe=teal!35!black,title={Before you move on}]
What does o peregrino mean? (\textquotedblleft{}The pilgrim\textquotedblright{}.) Say it once more.
\end{tcolorbox}
\section[a peregrinação]{a peregrinação — the pilgrimage}

\label{lesson:PT-C320-peregrinacao}



\begin{quote}
Before the new one: say the Portuguese for the ceremony, then the Portuguese for the pilgrim.
\end{quote}

\subsection*{You'll want to know: a peregrinação}
\textbf{a peregrinação} — \textquotedblleft{}the pilgrimage\textquotedblright{}.

\begin{quote}
Todos os anos, a minha tia faz uma peregrinação a pé.
\end{quote}

\begin{quote}
\emph{Every year my aunt goes on a pilgrimage on foot.}
\end{quote}

Built on \emph{peregrino}, ‘pilgrim’.

\begin{cousinweb}
\emph{peregrinação} is a learned borrowing from Latin \emph{peregrīnātiō} \textquotedblleft{}a travelling abroad\textquotedblright{}, from \emph{peregrīnus} \textquotedblleft{}foreign\textquotedblright{}, built on \emph{per} and \emph{ager} \textquotedblleft{}a field\textquotedblright{} (Lewis \& Short).
\end{cousinweb}

\begin{grammarlens}[title={What you've built}]
One more word for telling, explaining and describing at B1.
\end{grammarlens}

\subsection*{Your turn}
\begin{itemize}
\raggedright
  \item \emph{Say it:} \emph{a peregrinação}
  \item \emph{Say it:} \emph{a peregrinação}, once more
  \item \emph{Say it:} a sentence of your own with \emph{a peregrinação}
  \item \emph{Recall:} say the Portuguese for the ceremony, then the Portuguese for the pilgrim, then say \emph{a peregrinação} again
\end{itemize}

\begin{tcolorbox}[breakable,colback=teal!4,colframe=teal!35!black,title={Before you move on}]
What does a peregrinação mean? (\textquotedblleft{}The pilgrimage\textquotedblright{}.) Say it once more.
\end{tcolorbox}
